\documentclass[a4paper]{article}
\usepackage[top=0.9cm,bottom=0.86cm,left=1.47cm,right=1.47cm]{geometry}
\usepackage[T1]{fontenc}\usepackage{mathptmx}\usepackage{xcolor}\usepackage{enumitem}\usepackage[hidelinks]{hyperref}\usepackage{titlesec}
\pagestyle{empty}\setlength{\parindent}{0pt}\setlength{\parskip}{0pt}\setlength{\emergencystretch}{2em}\hyphenpenalty=10000\exhyphenpenalty=10000\tolerance=2000
\newcommand{\cvsection}[1]{\par\vspace{7.0pt}{\fontsize{11}{11}\selectfont\bfseries\MakeUppercase{#1}}\par\vspace{1pt}\hrule\vspace{2.5pt}}
\newcommand{\entry}[2]{{\bfseries #1} $|$ #2\par}
\newcommand{\skillrow}[2]{{\bfseries #1:} #2\par\vspace{1pt}}
\setlist[itemize]{leftmargin=12.95pt,labelsep=4pt,itemindent=0pt,topsep=0pt,before={\vspace{1.0pt}},itemsep=0.8pt,parsep=0pt,partopsep=0pt}
\begin{document}\fontsize{9.5}{9.9}\selectfont
\begin{center}{\fontsize{21.5}{22}\selectfont Sharan Teja Medikar}\par\vspace{3pt}{\fontsize{8.4}{9}\selectfont Guildford, UK | medikarsharanteja@gmail.com | +44 7824098307 | linkedin.com/in/sharanteja | github.com/sharantejamedikar | portfolio-mu-roan-71.vercel.app}\end{center}\vspace{2pt}
\colorbox{gray!14}{\parbox{\dimexpr\linewidth-2\fboxsep}{\textbf{SUMMARY}\par\vspace{2pt}Recently completed an MSc in Artificial Intelligence at the University of Surrey, with hands-on experience building AI/ML systems using Python, PyTorch and LLMs. Experienced across applied AI, machine learning, LLM evaluation, computer vision and software engineering, with a particular interest in building dependable AI systems for real-world problems. Seeking AI Engineering and Applied AI opportunities combining technical depth with end-to-end ownership. Currently have the right to work in the UK.}}
\par\vspace{1.0pt}\cvsection{Education}\entry{MSc in Artificial Intelligence}{University of Surrey | Sep 2025 - Sep 2026 | Result pending}\begin{itemize}\item \textbf{Modules:} Fundamentals of Machine Learning, AI and Sustainability, Ethics and Regulations of AI\item \textbf{Dissertation:} Self-Refining Code Generation Using Execution Feedback\end{itemize}\vspace{2.0pt}\entry{BTech in Computer Science and Engineering}{National Institute of Technology (NIT), Silchar | Jun 2021 - May 2025}\begin{itemize}\item \textbf{Modules:} Machine Learning, Data Structures and Algorithms\end{itemize}
\cvsection{Work Experience}\entry{Founder \& Technical Lead --- Police Night Patrolling System}{Government of Assam, Shri Bhumi District | Assam, India | Dec 2023 - Sep 2024}\begin{itemize}\item \textbf{Conceived and proposed} a real-time police patrolling platform to district government stakeholders, secured approval and assembled a five-person engineering team to deliver it from idea to operational handover.\item \textbf{Led development} of real-time tracking supporting 30 patrol vehicles and 60+ operational users, contributing to authentication, backend and real-time data functionality using Flutter, Firebase and Express.js.\item \textbf{Led testing, UAT and technical handover} for district-level operational deployment.\end{itemize}\vspace{2.5pt}\entry{Software Engineering Intern}{Fleckor Tech | Hyderabad, India | Jun 2024 - Jul 2024}\begin{itemize}\item \textbf{Developed frontend and backend components} contributing to 3+ core features including API integration and application functionality under sprint timelines.\item \textbf{Identified and resolved approximately 15-20 functional bugs} and supported deployment preparation and application testing.\end{itemize}
\cvsection{Projects}\entry{Self-Refining Code Generation Using Execution Feedback}{Python, PyTorch, Qwen2.5-Coder, CodeLlama}\begin{itemize}\item \textbf{Built an execution-feedback code-generation system} that generated solutions, executed tests and iteratively refined failures using structured runtime feedback.\item \textbf{Improved Qwen MBPP pass@1 from 71.9\% to 82.44\%} using adaptive execution-feedback refinement while retaining executable evaluation and bounded stopping.\item \textbf{Evaluated single-pass, best-of-5, adaptive and fixed-iteration refinement} across Qwen2.5-Coder and CodeLlama on HumanEval and MBPP, measuring reliability through executable outcomes.\item \textbf{Analysed convergence and failure modes} to compare when additional execution-feedback iterations improved or failed to improve generated solutions.\end{itemize}\vspace{2.5pt}\entry{AI Job Application Agent}{Python, Playwright, SQLite, Agentic Workflows}\begin{itemize}\item \textbf{Built an agentic job-search and application system} combining vacancy discovery, deterministic filtering, semantic matching, evidence-grounded preparation and browser-based ATS interaction.\item \textbf{Implemented persistent workflow state, duplicate prevention and validation gates} with human-in-the-loop controls for ambiguous or security-sensitive actions backed by automated tests.\item \textbf{Separated deterministic safety logic from semantic judgement} while preserving source provenance and auditable state transitions across the workflow.\end{itemize}\vspace{2.5pt}\entry{Vision-Language Model Prompt Learning and Adaptation}{CLIP, CoOp, CoCoOp, LoRA, Prompt Tuning}\begin{itemize}\item \textbf{Implemented and evaluated CLIP adaptation and prompt-learning methods} including CoOp, CoCoOp, KgCoOp, PromptSRC, ProGrad, LoRA, Tip-Adapter, CLIP-Adapter, MaPLe and fusion across EuroSAT and Flowers.\item \textbf{Improved EuroSAT from 53.99\% zero-shot accuracy to 89.70\% with Attention Fusion} while CoOp achieved 91.23\% Top-1 accuracy on Flowers across remote-sensing and fine-grained visual domains.\item \textbf{Compared adaptation baselines on EuroSAT} with LoRA achieving 64.84\%, CoCoOp 83.14\%, ProGrad 84.99\%, KgCoOp 85.59\% and Fusion MLP 88.98\% accuracy.\end{itemize}\vspace{2.5pt}\entry{Physics-Informed Environmental ML}{Python, PyTorch, XGBoost, LightGBM}\begin{itemize}\item \textbf{Developed a physics-informed LSTM over 28,632 site-year observations across 26,446 locations} achieving RMSE 32.053 and R$^2$ 0.291 on temporal 2024 evaluation.\item \textbf{Improved RMSE by 34.5\% over XGBoost and 9.7\% over a standard LSTM} while evaluating compression, interpretability and environmental-equity behaviour.\item \textbf{Reduced model size with float16 quantisation} from approximately 1.8 MB to 908 KB with essentially no RMSE degradation.\end{itemize}
\cvsection{Skills}\skillrow{Programming \& Data}{Python, SQL, JavaScript, NumPy, Pandas}\skillrow{AI \& Machine Learning}{PyTorch, scikit-learn, Computer Vision, LLM Evaluation, CLIP}\skillrow{Generative AI}{Qwen, CodeLlama, AI Agents, Prompt Engineering}\skillrow{Engineering \& Tools}{APIs, Testing, Git, Linux, Firebase}
\end{document}